\documentclass[10pt,a4paper]{article}
\usepackage[utf8]{inputenc}
\usepackage[T1]{fontenc}
\usepackage{amsmath,amssymb}
\usepackage{graphicx}
\usepackage{booktabs}
\usepackage{hyperref}
\usepackage{xcolor}
\usepackage[margin=2.5cm]{geometry}

\title{RAMP: Data-Free Mixed-Precision Quantization\\for Hybrid SSM-MoE Models on Consumer Hardware}

\author{
  K\'evin Remondi\`ere\\
  Independent Researcher, Oloron-Sainte-Marie, France\\
  \texttt{kevin.remondiere@gmail.com}\\
  ORCID: \href{https://orcid.org/0009-0008-2443-7166}{0009-0008-2443-7166}\\
  \url{https://github.com/AIdevsmartdata/ramp-quant}
}

\date{March 2026}

\begin{document}
\maketitle

\begin{abstract}
We present RAMP (RL-guided Adaptive Mixed-Precision), a quantization pipeline that produces hardware-optimized GGUF files for hybrid SSM-MoE models without requiring calibration data during the allocation search. Applied to Qwen3.5-35B-A3B (30~GDN + 10~attention layers, 256~experts), RAMP~v2 produces a 15.2\,GB file (3.78~BPW) running at 90~tok/s on RTX~5060~Ti 16\,GB with 30/30 on functional benchmarks. The pipeline combines NSDS data-free sensitivity scoring, proxy-based loss estimation, and evolutionary search under VRAM constraints. We document the complete engineering arc including 7~failed builds, a QuaRot dead end (PPL~49{,}524), and the discovery that rotation-based methods are fundamentally incompatible with GDN recurrent layers. Code, model, and imatrix released.\footnote{\url{https://github.com/AIdevsmartdata/ramp-quant}, \url{https://huggingface.co/Kevletesteur/Qwen3.5-35B-A3B-RAMP-v2-15G}}
\end{abstract}

\section{Introduction}

Quantizing hybrid SSM-MoE models differs fundamentally from quantizing dense transformers. Qwen3.5-35B-A3B has three structurally distinct components with different sensitivity profiles:

\begin{itemize}
  \item \textbf{GDN recurrent layers} (30/40): errors propagate multiplicatively through the state recurrence $M_t = \alpha_t M_{t-1}(I - k_t q_t^\top) + k_t v_t^\top$. At timestep $t$, quantization error grows as $O(e^t)$.
  \item \textbf{Attention layers} (10/40): standard GQA with 16~Q-heads (dim=512) and 2~KV-heads (dim=256). Tolerant to 4--5~bit quantization.
  \item \textbf{MoE experts} (256 per layer, top-8): represent 93\% of parameters but are sparsely activated. Highly compressible.
\end{itemize}

Standard uniform quantization (e.g., IQ3\_S everywhere) ignores this structure. We show that per-tensor mixed-precision allocation, guided by data-free sensitivity analysis, produces significantly better quality-size tradeoffs.

\section{The Road to RAMP: Failed Approaches}

Before arriving at RAMP, we attempted four rotation-based quantization methods. All failed on this architecture.

\subsection{QuaRot: Catastrophic Failure}

QuaRot~\cite{spinquant2024} was applied shard-by-shard (14~safetensors) to the BF16 model. Result: \textbf{complete garbage} at inference. The SSM $\alpha$ (decay gate) was shifted by $105{\times}$. Root cause: Hadamard rotation $R$ applied to weight $W$ changes the output distribution. The residual connection adds the \emph{unrotated} input to the \emph{rotated} output. RMSNorm does not commute with arbitrary orthogonal rotations. For recurrent layers, the rotated state accumulates non-linearly --- $R$ cannot be absorbed.

\subsection{OptRot: PPL = 49{,}524}

OptRot~\cite{optrot2024} applies data-free Givens rotations minimizing $\|RW\|^4$ (kurtosis). Applied to 40~attention tensors. Kurtosis reduced by 25\% on average, but \textbf{PPL = 49{,}524} (vs.\ $\sim$3.4 baseline). Same root cause: per-matrix rotation without cross-layer absorption.

\subsection{ParoQuant: OOM}

ParoQuant~\cite{paroquant2024} handles absorption correctly but requires the full model in memory. On RTX~5060~Ti + 32\,GB RAM: 256~MoE experts $\to$ \texttt{grouped\_mm} exceeds 16\,GB VRAM at batch\_size=1. Adding 40\,GB swap allowed model loading but OOM on forward pass. \textbf{Requires A100 80\,GB.}

\subsection{EvoPress: Same OOM}

EvoPress~\cite{evopress2024} uses KL-divergence fitness requiring full-model forward. Same 70\,GB RAM barrier.

\textbf{Lesson:} On consumer hardware (16\,GB VRAM, 32\,GB RAM), rotation and calibration-based methods are infeasible for 35B MoE models. Data-free approaches are the only option.

\section{RAMP Pipeline}

RAMP operates in 6~phases without loading the full model:

\subsection{Phase 1: GGUF Analysis}

Parse the GGUF file structure: 733~tensors, 40~layers, 263~decision groups. Classify each tensor by role: expert (120~groups, 32.2B elements, 93\%), SSM (30~groups, 256.6M), attention (40~groups, 272.7M), norm, gate.

\subsection{Phase 2: NSDS Sensitivity (Data-Free)}

NSDS (Normalized Data-free Sensitivity)~\cite{nsds2024} scores each tensor group using kurtosis and SVD rank, without any calibration data. Computation: $\sim$5~minutes, processes tensors one-at-a-time (peak RAM = 1~tensor).

Results confirm the GDN sensitivity hierarchy:
\begin{itemize}
  \item Norms: NSDS = 1.000 (highest)
  \item SSM gates ($\alpha$, $\beta$): NSDS = 0.95+
  \item Attention K/Q: NSDS = 0.85--0.95
  \item Expert FFN (mid layers 17--23): NSDS = 0.748 (lowest)
\end{itemize}

\subsection{Phase 3: Proxy Model}

Build a lookup table of quantization error per tensor per type (8~types: IQ2\_XXS to Q8\_0). Each entry: round-trip error $\|W - \text{dequant}(\text{quant}(W))\|$. Combined with NSDS sensitivity, this gives an instant proxy loss for any allocation.

\subsection{Phase 4: Search}

Two search strategies under a VRAM budget constraint:

\textbf{Greedy (ScaleBITS-style):} 757~steps in 7.89\,s. Iteratively upgrades the most sensitive tensor that fits the budget. Result: score=4.765, 14.71\,GB.

\textbf{Evolutionary (128~pop, 200~gen):} 25{,}600 evaluations in 603\,s. Converges at generation $\sim$120 (stagnation). Result: score=5.201. \textbf{Greedy wins} --- the budget is too tight for evolutionary diversity to help.

\subsection{Phase 5: Validate}

Verify the allocation fits hardware: VRAM with KV cache, ncmoe offloading headroom.

\subsection{Phase 6: Build}

Generate a \texttt{llama-quantize --custom-q} command with regex-based per-tensor overrides. The actual quantization uses stock \texttt{llama-quantize} --- no custom code needed at this stage.

\section{Results}

\subsection{RAMP v1 (March 29, Morning)}

\begin{table}[h]
\centering
\caption{RAMP v1 vs.\ baselines.}
\begin{tabular}{lccccc}
\toprule
\textbf{Model} & \textbf{Size} & \textbf{BPW} & \textbf{tok/s} & \textbf{Bench} & \textbf{Notes} \\
\midrule
IQ3\_S Unsloth UD & 13.0\,GB & 3.44 & 69 & --- & Stock, 3 slots \\
IQ3\_S custom-mix & 14.71\,GB & 3.56 & 98.6 & 20/20 & 317 manual overrides \\
Q4\_K\_M & 21.0\,GB & $\sim$4.5 & 62.3 & --- & OOM at 128K \\
\textbf{RAMP v1} & \textbf{16.6\,GB} & \textbf{4.02} & \textbf{92.4} & --- & Q3\_K\_M base, no imatrix \\
\bottomrule
\end{tabular}
\end{table}

RAMP v1 used Q3\_K\_M as base (no imatrix required) with Q5\_K/Q8\_0 overrides on critical paths. Two commands, zero custom code, 10~minutes total.

\subsection{RAMP v2 (March 29, Evening)}

Key improvement: generate a custom imatrix from RAMP~v1 itself (3.5~min, 200~chunks, domain-calibrated corpus: 100~kine FR, 50~code, 30~math, 30~tool-calling, 68~training pairs). Then requantize with IQ3\_S base.

\begin{table}[h]
\centering
\caption{RAMP v2 vs.\ v1 improvements.}
\begin{tabular}{lcc}
\toprule
& \textbf{RAMP v1} & \textbf{RAMP v2} \\
\midrule
Size & 16.6\,GB (4.02 BPW) & \textbf{15.2\,GB (3.78 BPW)} \\
Base type & Q3\_K\_M & IQ3\_S + imatrix \\
ssm\_out & Q5\_K & \textbf{Q8\_0} \\
attn\_k & Q5\_K & \textbf{Q6\_K} \\
token\_embd & Q5\_K & \textbf{Q6\_K} \\
Gen speed & 92.4 tok/s & 90.3 tok/s \\
Functional bench & good & \textbf{30/30} \\
\bottomrule
\end{tabular}
\end{table}

The size reduction (16.6 $\to$ 15.2\,GB) comes from IQ3\_S ($\sim$3.44 BPW) vs.\ Q3\_K\_M ($\sim$3.56 BPW) as base. The quality improvement comes from promoting ssm\_out to Q8\_0 (critical in GDN sensitivity hierarchy) and using a domain-calibrated imatrix.

\subsection{The Night of Builds}

Table~\ref{tab:builds} shows the 7~build attempts on March 28--29.

\begin{table}[h]
\centering
\caption{Build attempts chronology.}
\small
\begin{tabular}{llll}
\toprule
\textbf{\#} & \textbf{Source} & \textbf{Result} & \textbf{Failure mode} \\
\midrule
1 & BF16 QuaRot & Garbage & SSM corrupted by rotation \\
2 & BF16 QuaRot + fallback & Garbage & Rotated experts contaminate \\
3 & convert\_hf $\to$ F16 & 137\,GB file & F32 experts (2\,GB each) \\
4 & Requant from build~1 & Garbage & Source was garbage \\
5 & convert\_hf $\to$ BF16 & Correct (66\,GB) & Slow (redundant with Unsloth) \\
6 & BF16 clean + custom-q & \textbf{17\,GB, works} & First success \\
7 & BF16 rote + hybrid & Garbage & Any rotation = garbage \\
\bottomrule
\end{tabular}
\label{tab:builds}
\end{table}

\textbf{Root cause:} 4/7 failures from trying to use QuaRot-rotated weights. The epiphany: Unsloth already provides correct BF16 GGUF with all SSM transformations pre-applied.

\section{MTP GGUF: Parallel Story}

We also built tools to inject Multi-Token Prediction tensors into GGUF files (4~script versions over 2~weeks). The resulting MTP GGUF (18.4\,GB, 753~tensors) achieved 84.8\% acceptance rate but \textbf{slowed inference from 93 to 47~tok/s} because the MTP layer is itself MoE (256~experts on CPU). The MTP was pruned for production. Scripts and model released.\footnote{\url{https://huggingface.co/Kevletesteur/Qwen3.5-35B-A3B-IQ3_S-MTP}}

\section{Lessons Learned}

\begin{enumerate}
  \item \textbf{Rotation methods are incompatible with GDN.} QuaRot, OptRot, and any Hadamard/Givens rotation without full cross-layer absorption corrupt recurrent state. MambaQuant~\cite{mambaquant2024} confirms: 21\% accuracy drop even at W8A8.
  \item \textbf{MoE experts are the most compressible.} 93\% of parameters, sparsely activated. Couches 17--23 tolerate IQ3\_XXS; periphery (0, 38--39) needs Q4\_K.
  \item \textbf{Don't write custom code when the tool exists.} Our 900-line \texttt{ramp\_quantize.c} is obsolete --- \texttt{llama-quantize --custom-q} does the same thing.
  \item \textbf{Always start from Unsloth BF16.} SSM transformations (deinterleave, $-\exp$, norm+1) pre-applied. \texttt{convert\_hf\_to\_gguf.py} produces 2$\times$ larger files.
  \item \textbf{Consumer hardware = data-free methods only.} ParoQuant, EvoPress, SpinQuant all OOM. NSDS + proxy search works locally in minutes.
  \item \textbf{The imatrix bootstrapping trick:} Use RAMP~v1 (no imatrix, Q3\_K\_M) to generate a domain-calibrated imatrix, then requantize to IQ3\_S for RAMP~v2.
\end{enumerate}

\section{Related Work}

GPTQ~\cite{gptq2023} and AWQ~\cite{awq2023} require calibration data and full model in memory. SpinQuant~\cite{spinquant2024} learns rotations but needs the full model. EvoPress~\cite{evopress2024} uses evolutionary search with KL fitness (same OOM). NSDS~\cite{nsds2024} provides data-free sensitivity --- we build on this. Unsloth Dynamic~2.0 is the production baseline. RAMP~\cite{ramp2024} uses RL for allocation --- we replace RL with proxy + greedy/evolutionary search for consumer feasibility.

\section{Conclusion}

RAMP demonstrates that data-free mixed-precision quantization can produce production-quality GGUF files for hybrid SSM-MoE models on consumer hardware. The key insight is that GDN sensitivity hierarchy (gates $>$ attention $>$ shared experts $>$ routed experts) enables aggressive compression of 93\% of parameters while protecting the 7\% that matter. RAMP~v2 runs at 90~tok/s with 30/30 functional accuracy on RTX~5060~Ti --- comparable to Q4\_K\_M quality at 71\% of the size.

\begin{thebibliography}{10}

\bibitem{gptq2023} Frantar, E., et~al. GPTQ: Accurate post-training quantization for generative pre-trained transformers. \emph{ICLR}, 2023.

\bibitem{awq2023} Lin, J., et~al. AWQ: Activation-aware weight quantization for LLM compression and acceleration. \emph{MLSys}, 2024.

\bibitem{spinquant2024} Liu, Z., et~al. SpinQuant: LLM quantization with learned rotations. \emph{arXiv:2405.16406}, 2024.

\bibitem{evopress2024} IST-DASLab. EvoPress: Evolutionary search for GGUF quantization. \emph{ICML}, 2025.

\bibitem{nsds2024} NSDS: Normalized data-free sensitivity for mixed-precision quantization. \emph{arXiv:2603.17354}, 2024.

\bibitem{ramp2024} RAMP: RL-guided adaptive mixed-precision quantization. \emph{arXiv:2603.17891}, 2024.

\bibitem{optrot2024} OptRot: Optimal rotations for quantization. \emph{NeurIPS}, 2025.

\bibitem{paroquant2024} ParoQuant: Pairwise rotations for quantization. \emph{ICLR}, 2026.

\bibitem{mambaquant2024} MambaQuant: Quantizing recurrent state space models. \emph{ICLR}, 2025.

\end{thebibliography}

\end{document}
